% Build-time font warm-up — NOT part of the CV. Never compiled into an artifact.
%
% cv/Dockerfile compiles this once so the EC/LModern fonts and LaTeX format
% files are baked into the toolchain image. That turns the per-run cost of the
% PDF from ~2.0 s (regenerating fonts via mktexpk every time) into ~0.44 s.
%
% The preamble deliberately mirrors cv/cv.tex's essentials, because the point is
% to prime the same engine configuration and font encoding. The body then uses
% every font family and every size the CV needs, so their PK files all exist
% before a real run.
%
% If cv.tex later starts using a font shape or size that is not exercised here,
% the first run will regenerate just that font (~1.5 s, once) — the build still
% succeeds, so this file is a performance hint, not a correctness dependency.

% The preamble and class options MUST mirror cv/cv.tex exactly, because EC font
% sizes are derived from the document's base point size. A bare \documentclass
% {article} is 10pt, so \normalsize is ecrm1000 — but cv.tex is 11pt, where
% \normalsize is 10.95pt (ecrm1095). Warming up 1000-series fonts against an
% 11pt document bakes 16 fonts that are never used and leaves the 8 that ARE
% used to be generated at run time, so the warm-up silently does nothing.
% (Measured: the real CV generated ecrm1095, ecbx1095, ecbi1095, ecti1095,
% ecsi1095, ectt1095, tcrm1095, ecbx1440 on every run.)
%
% If cv.tex's class options or fontencoding change, update this file to match or
% the warm-up becomes a no-op again.

\documentclass[11pt,letterpaper]{article}

\usepackage[T1]{fontenc}
\usepackage[utf8]{inputenc}
\usepackage[margin=1in]{geometry}
\usepackage[hidelinks]{hyperref}
\usepackage[expansion=false]{microtype}
\DisableLigatures{encoding = T1, family = *}

\setlength{\parindent}{0pt}
\setlength{\parskip}{6pt}
\pagestyle{empty}

\begin{document}

% Systematically exercise EVERY EC family at EVERY size the document class can
% produce, rather than hand-listing the combinations that happened to be needed
% today. Hand-listing is what broke the first attempt: it missed ecbx1440, so one
% font was still generated at run time. A sweep cannot go stale when cv.tex gains
% a new heading style, and the cost is a few hundred KB of PK files, once.
%
% \usefont takes NFSS family names, which are the CM names (cmr, cmss, cmtt) —
% NOT the EC file names (ecrm, ecss). The .fd files map NFSS -> EC file, e.g.
% t1cmr.fd has \EC@family{T1}{cmr}{bx}{it}{ecbi}, so {cmr}{bx}{it} produces
% ecbi1095. Using ecrm/ecss here silently selects nothing and bakes no fonts —
% it looks right and does nothing, which is exactly how the first sweep failed.
% The eight EC families are the cross product of cmr/cmss/cmtt with series/shape.
\newcommand{\sweep}[1]{%
  \fontsize{#1}{#1}\selectfont
  \usefont{T1}{cmr}{m}{n}rm \usefont{T1}{cmr}{bx}{n}bx \usefont{T1}{cmr}{m}{it}it
  \usefont{T1}{cmr}{bx}{it}bi \usefont{T1}{cmr}{m}{sc}sc \usefont{T1}{cmr}{m}{sl}sl
  \usefont{T1}{cmss}{m}{n}ss \usefont{T1}{cmss}{bx}{n}sx \usefont{T1}{cmss}{m}{sl}ssi
  \usefont{T1}{cmtt}{m}{n}tt\par
}

% TS1 (the "text companion" encoding) supplies symbols such as \textbullet, which
% cv.tex uses as the separator on its contact line. Its roman font is tcrm, and it
% is selected via a DIFFERENT encoding, so the T1 sweep above never reaches it.
% Measured: without this, tcrm1095 was generated on every single run.
\newcommand{\sweepTSone}[1]{%
  \fontsize{#1}{#1}\selectfont
  \usefont{TS1}{cmr}{m}{n}tcrm \usefont{TS1}{cmr}{bx}{n}tcbx
  \usefont{TS1}{cmss}{m}{n}tcsi \usefont{TS1}{cmtt}{m}{n}tctt\par
}

% The sizes the 11pt article class produces, from \tiny through \Huge.
\sweep{6}      % \tiny
\sweep{8}      % \scriptsize
\sweep{9}      % \footnotesize
\sweep{10}     % \small
\sweep{10.95}  % \normalsize
\sweep{12}     % \large
\sweep{14.4}   % \Large
\sweep{17.28}  % \LARGE
\sweep{20.74}  % \huge
\sweep{24.88}  % \Huge

% The same sweep for TS1 symbols (see \sweepTSone above).
\sweepTSone{10.95}
\sweepTSone{14.4}

% \textbullet is how cv.tex separates contact details, so exercise the exact
% markup path too, not just the font selection. Plain text rather than cv.tex's
% \ph{} macro, which is deliberately not defined here.
\usefont{T1}{cmr}{m}{n}\selectfont
EMAIL \textbullet\ PHONE \textbullet\ CITY\\

% Also run the normal markup paths, so anything the sweep does not cover directly
% (list setup, bold+italic nesting, the sans-serif \cvwhen style) is exercised.
\usefont{T1}{ecrm}{m}{n}\selectfont
\textbf{Title} \textsf{Date} \textit{Placeholder}
\textbf{\textsf{Bold Sans}} \textit{\textsf{Italic Sans}}
\textbf{\textit{Bold Italic}} \textbf{\itshape Bold Italic Two}

\begin{itemize}
  \item A bullet so the itemize font setup is cached.
\end{itemize}

Alden • github.com/alden-g878 • Universit\'e na\"ive — accents and separators.

\end{document}
